%==============================================================================
% Experiments, Results, and Analysis
%
% Drop-in body for Section 5 (Experiments). The enclosing \section{Experiments}
% and \label{sec:experiments} live in main.tex; this file supplies the
% condensed body subsections. All numbers are the verified means/per-category
% scores read off results/<system>/helix-{small,medium}.json and cross-checked
% against the ground-truth anchors. Public field = four systems (gbrain,
% zep-cloud, mem0-platform, graphify-oss); the only medium-tier run completed to
% date is gbrain. Per-system autopsies, the gbrain mechanism subsection, the full
% readiness case write-ups, per-run dollar figures, and judge token detail are in
% Appendix~\ref{app:extended} (Extended Results).
%==============================================================================

\subsection{Metrics and Setup}
\label{sec:exp:metrics}

\bench{} scores each answer with a \emph{rubric-weighted} judgment in $[0,1]$
rather than exact match. Every question ships with a weighted rubric (Stage~E,
\S\ref{sec:design:questions}): a small set of sub-criteria---typically the
specific fact, the date or time qualifier, the supersession or contradiction
relation, and the justification chain---each with an explicit weight. An answer
earns the weighted fraction of sub-criteria it satisfies, so partial credit
reflects \emph{which} part of an organizationally grounded answer a system got
right. The headline metric is the mean of this rubric score over the questions
in a tier. We report exact match only as a secondary diagnostic: it is
effectively zero across the board, because correct answers are multi-clause
statements no system reproduces verbatim, which is exactly why a sub-point rubric
is the appropriate instrument. An LLM judge with extended thinking scores each
answer against the per-sub-criterion rubric, awarding each sub-point
independently before summing, under the same configuration for every system
(token budgets, temperature, and output caps in Appendix~\ref{app:extended}).

The judge also labels each non-perfect answer with one of five mutually exclusive
failure modes---\texttt{retrieval\_miss}, \texttt{hallucination},
\texttt{wrong\_synthesis}, \texttt{wrong\_abstain}, and \texttt{none}---which lets
us separate \emph{retrieval} failures from \emph{reasoning} failures, a
distinction the scaling analysis below turns on.

\paragraph{Setup and answerer policy.}\label{sec:exp:setup}
Each system is evaluated in isolation: the tier corpus is ingested into a
\emph{fresh} store (scope isolated per tier), the system is driven through a
readiness gate (\S\ref{sec:exp:readiness}), and the tier's questions are then
issued and scored; hosted systems run over their public APIs and self-hosted
systems in pinned Docker images (pins in \S\ref{sec:baselines}/Appendix~\ref{app:hyperparams}).
Retrieval and answer synthesis are decoupled: a system that ships its own
synthesizer is run end-to-end with it (\texttt{gbrain} via native \texttt{think}),
the artifact a user would deploy, while retrieval-only systems
(\texttt{zep-cloud}, \texttt{mem0-platform}, \texttt{graphify-oss}) share one
fixed neutral answerer (Claude Sonnet~4.6) so differences reflect the memory
layer; top-$K=10$ wherever configurable.

\paragraph{Cost.}
End-to-end cost is $\approx$\$0.02--\$0.04 per question, dominated by the judge
and answerer; \texttt{graphify-oss}'s server-side Sonnet extraction dominates its
all-in cost ($\approx$\$0.13 per question at medium), with the full breakdown and
the \texttt{cost\_per\_query\_usd} metering caveat in
Appendix~\ref{app:extended}.

\paragraph{Statistical caveat (stated up front).}
These are \textbf{single-seed} results over \textbf{small} question sets (11
small, 73 medium), so we attach no confidence intervals and treat sub-$0.1$
between-system gaps as indicative rather than significant. The full treatment of
this variance---including the replication numbers---is in \S\ref{sec:limitations},
the body home for those figures.

\subsection{Main Results}
\label{sec:exp:main}

Table~\ref{tab:main-results} reports the headline rubric means. The central
finding is not that one system wins: it is that \textbf{organizational-scale
memory is unsolved}. The strongest currently-available public system,
\texttt{gbrain}, reaches only $0.394$ on the small tier and $0.428$ on the
medium tier; across both tiers the rest of the public field clusters between
$0.142$ and $0.303$. There is large headroom on every tier and in every
category, and the spread between the best public system and the rest, while
clear, is itself modest relative to the gap to a competent answer. We
deliberately frame \texttt{gbrain} as the strongest \emph{currently-available
public} system rather than as a solution: a $0.39$ ceiling on a benchmark whose
questions a human analyst with access to the corpus could largely answer is the
paper's motivating result, not a leaderboard victory. ``Leads'' is not
``solves'': \texttt{gbrain} scores $0.000$ on small-tier contradiction (C6) and
$0.122$ on bi-temporal (C3), and its $0.394$/$0.428$ means leave roughly
six-tenths of the achievable rubric credit on the table. The headline of the
benchmark is this remaining headroom, not the identity of the front-runner.

\begin{table}[t]
  \centering
  \caption{\textbf{End-to-end results on \bench{}: mean rubric score and
    per-category breakdown (higher is better).} A single consolidated view of
    both tiers. Top panel: the four public memory systems on the small tier
    (11 questions, 121 artifacts). Bottom panel: the same four on the medium tier
    (73 questions, 443 artifacts). Category columns abbreviate the
    six \bench{} categories (full definitions in Table~\ref{tab:categories}):
    Supers.~=~supersession (C1), Prov.~=~decision provenance (C2),
    Bi-temp.~=~bi-temporal as-of (C3), Audit~=~audit replay (C4),
    Justif.~=~justification chain (C5), Contra.~=~contradiction (C6). Best
    value per column within each tier in bold. Single seed; see \S\ref{sec:exp:setup}
    for the variance caveat. A controlled-extractor probe (\texttt{mem0} OSS +
    Sonnet, not a shipped product; Appendix~\ref{app:extended}) scores a $0.479$
    small-tier mean.}
  \label{tab:main-results}
  \begin{tabular}{l S[table-format=1.3]
      S[table-format=1.3] S[table-format=1.3] S[table-format=1.3]
      S[table-format=1.3] S[table-format=1.3] S[table-format=1.3]}
    \toprule
    {System} & {Mean} & {Supers.} & {Prov.} & {Bi-temp.} & {Audit} & {Justif.} & {Contra.} \\
    \midrule
    \multicolumn{8}{l}{\emph{Small tier} --- 11 questions, 121 artifacts ($\approx$60K tokens)} \\
    \texttt{gbrain}        & \bfseries 0.394 & \bfseries 0.675 & \bfseries 0.762 & 0.122 & \bfseries 0.533 & \bfseries 0.411 & 0.000 \\
    \texttt{zep-cloud}     & 0.252 & 0.525 & 0.625 & 0.172 & 0.000 & 0.117 & \bfseries 0.133 \\
    \texttt{mem0-platform} & 0.221 & 0.000 & 0.400 & \bfseries 0.317 & 0.000 & 0.183 & \bfseries 0.133 \\
    \texttt{graphify-oss}  & 0.206 & 0.362 & 0.250 & 0.028 & 0.067 & 0.406 & 0.033 \\
    \midrule
    \multicolumn{8}{l}{\emph{Medium tier} --- 73 questions, 443 artifacts ($\approx$200K tokens)} \\
    \texttt{gbrain}        & \bfseries 0.428 & \bfseries 0.540 & \bfseries 0.612 & \bfseries 0.307 & \bfseries 0.489 & \bfseries 0.216 & 0.214 \\
    \texttt{mem0-platform} & 0.303 & 0.403 & 0.492 & 0.000 & 0.348 & 0.059 & \bfseries 0.431 \\
    \texttt{zep-cloud}     & 0.213 & 0.252 & 0.271 & 0.026 & 0.402 & 0.020 & 0.194 \\
    \texttt{graphify-oss}  & 0.142 & 0.068 & 0.185 & 0.067 & 0.264 & 0.141 & 0.000 \\
    \bottomrule
  \end{tabular}
\end{table}

\paragraph{Reading the table.}
Two orderings emerge: small tier \texttt{gbrain}~$\gg$ \texttt{zep-cloud}~$>$
\texttt{mem0-platform}~$>$ \texttt{graphify-oss}; medium tier \texttt{gbrain}~$\gg$
\texttt{mem0-platform}~$>$ \texttt{zep-cloud}~$>$ \texttt{graphify-oss}. The gap
between \texttt{gbrain} and the rest is the only clearly significant separation at
both scales; the lower-ranked systems are closer together, so we treat fine
separations among them cautiously, and exact match is near zero everywhere,
confirming the rubric mean is the informative signal. Notably the order among the
trailing three is \emph{not} preserved across tiers---\texttt{mem0-platform} and
\texttt{zep-cloud} swap---which we return to in \S\ref{sec:exp:scaling}. A
controlled \texttt{mem0} OSS probe (extractor swapped to Sonnet~4.6) localizes
that system's bottleneck to its undisclosed extractor, with full framing in
Appendix~\ref{app:extended}.

\subsection{Analysis by Category}
\label{sec:exp:bytype}

The per-category panels of Table~\ref{tab:main-results} expose where the
difficulty concentrates, and the hardest categories are exactly the ones
\bench{} was built to probe. \textbf{Contradiction (C6) and bi-temporal (C3) are
cliffs already on the small tier}, near zero or below $0.32$ across the field.
\textbf{Justification chains (C5) hold up better at small scale but collapse as
the corpus grows.} The categories the benchmark targets stay near the floor as
the corpus grows---at medium scale no system clears $0.31$ on C3 or $0.22$ on
C5---so fine-grained entity, event, and as-of reasoning is precisely the behavior
no public system delivers, and it does not improve with more data. The
per-category profile also explains where each non-\texttt{gbrain} system lands and
why \texttt{gbrain} leads; the per-system failure-mode autopsies and the
\texttt{gbrain} retrieval/synthesis mechanism are in Appendix~\ref{app:extended}.

\subsection{Scaling from Small to Medium}
\label{sec:exp:scaling}

The medium tier triples the corpus ($443$ artifacts) and roughly quintuples the
question count ($73$). Running all four systems on both tiers shows that
\textbf{degradation with scale is system-specific, not uniform}, and the
small-tier ordering does \emph{not} simply persist: \texttt{gbrain} is
scale-robust ($0.394 \rightarrow 0.428$), \texttt{mem0-platform} improves
($0.221 \rightarrow 0.303$), \texttt{zep-cloud} degrades ($0.252 \rightarrow
0.213$), and \texttt{graphify-oss} degrades most in relative terms ($0.206
\rightarrow 0.142$). The net effect is that \texttt{mem0-platform} and
\texttt{zep-cloud} \emph{swap} ranks between tiers while \texttt{gbrain} extends
its lead. This is a stronger result than ``the ranking holds'': an architecture's
behavior at conversational small scale does not predict its behavior at
organizational scale, which is precisely the regime \bench{} is built to expose.
The per-system causal one-liners, edge counts, and \texttt{gbrain}'s medium-tier
failure-mode shift are in Appendix~\ref{app:extended}.

\subsection{Readiness Gating: A Methodological Result}
\label{sec:exp:readiness}

A fair comparison of \emph{hosted, asynchronous} memory systems requires a
\textbf{readiness gate} before querying, and getting this wrong materially
distorts the leaderboard. Synchronous systems have finished ingesting when the
ingest call returns; the two hosted async systems do not---extraction continues
server-side, so an ``ingest returned'' signal or a naive count-plateau drain can
declare readiness over an empty or half-built store. We treat the design of a
correct drain as a contribution in its own right: hosted-memory benchmarking must
(i)~gate on a \emph{probe search} that returns non-zero, stable results, not on an
ingest-returned signal or a count plateau, and (ii)~report per-category, not only
aggregate, scores---because graph completeness can move recall and reasoning in
opposite directions under a flat mean. Both points have leaderboard-changing
force: under the corrected gate \texttt{mem0-platform} rises $0.141 \rightarrow
0.221$, while a proper \texttt{zep-cloud} settle barely moves the tier mean yet
swings per-category scores massively. These fixes were necessary to produce the
numbers in Table~\ref{tab:main-results} and are, in our view, prerequisites for
\emph{any} fair evaluation of asynchronous hosted memory, not artifacts of this
benchmark. The two full case write-ups (the \texttt{mem0} drain mechanism and the
\texttt{zep} half-built-graph deltas with their \texttt{Retry-After}/404-race
plumbing) are in Appendix~\ref{app:extended}.
